\chapter{工程图（TechDraw）：从 3D 到二维蓝图}

工程图工作台把 3D 模型变成可交付的二维图纸：视图、剖视、局部放大、尺寸、符号，最后导出 PDF/SVG/DXF。
它是全流程（建模 → 装配 → 出图）的最后一站。

\section{页面与模板}

\begin{enumerate}
  \item 底部切到 \ui{工程图}；空状态显示 \ui{还没有工程图} 与 \ui{新建工程图页}；
  \item 点 \ui{新建工程图页}（命令 \cmd{TechDraw\_PageDefault}）；
  \item 页面使用偏好设置里的默认模板；没有设置时使用 A4 横向的 \file{Default\_Template\_A4\_Landscape.svg}；
  \item 模板字段（作者、日期、标题、比例、图号）可以被自动填充（取文档的 CreatedBy/Company/Label），
        也可以点字段直接改值。
\end{enumerate}

页面对象出现在特征树里；点页面行切换到它。双击视图打开属性（比例、方向、X/Y 位置等）；
\ui{重绘页面} 刷新所有视图。

\section{视图类型}

\begin{center}
\small
\begin{tabularx}{\linewidth}{@{}lX@{}}
\toprule
\textbf{视图} & \textbf{操作} \\
\midrule
视图 View & 选零件（可多选）→ 选方向（前/上/右/后/下/左/等轴测）→ 放在页面中心 \\
投影组 ProjectionGroup & 选零件 → 多选投影方向（默认前/上/右）→ 自动排成标准的第三角视图 \\
剖视图 SectionView & 先选基视图 → 选方向（右/左/上/下）→ 过视图中心剖切；符号自动 A/B/C，放在右（或下）侧 60 mm \\
复杂剖视 & 选基视图 → 选切割草图/线框 → 偏移策略 \\
局部详图 DetailView & 选基视图 → 在页面上点圆心 → 半径（默认 10）→ 自动生成 2× 放大、字母编号，放在 +80/+40 mm \\
活动视图 ActiveView & 当前 3D 视口的截图（宽 100 mm）\\
图像 Image & 插入图片（固定 100 × 100 mm）\\
符号 / 粗糙度 & 从资源 SVG 里选，插在页面中心 \\
注释 / 富文本 & 先输入文字，再点位置 \\
气球 Balloon & 选视图 → 点指向点 → 点气泡位置，编号自动 \\
引线 Leader & 选视图 → 点箭头点 → 继续点 → 完成 \\
焊接符号 & 选引线 → 选类型（当前版本类型参数未生效）\\
裁剪组 ClipGroup & 框一个区域并加/减视图（当前版本不实际裁剪）\\
\bottomrule
\end{tabularx}
\end{center}

还有：\ui{移动视图}、\ui{共享视图}（同一视图放多页）、\ui{叠放}、\ui{锁定}。

\section{尺寸标注}

在页面上点视图的边/顶点/面，再点尺寸命令。

\subsection{智能尺寸}

\ui{尺寸} 会按选择自动判断：

\begin{itemize}
  \item 选一个面 → 面积；
  \item 选两条边 → 角度；
  \item 选三个顶点 → 三点角；
  \item 选圆/圆弧边 → 询问直径还是半径；
  \item 其它 → 询问长度/水平/垂直。
\end{itemize}

\subsection{专用尺寸}

长度、水平、垂直、角度、三点角、面积、半径、直径、弧长、范围……每个命令都有选择校验，选错会提示
\ui{Incorrect selection: select …}。另有链式、坐标、倒角尺寸，以及基于所选尺寸的“对齐/层叠”。

尺寸会自动放置（一般偏 7 mm、半径 1.1r、角度 15 mm），\textbf{没有“点一下放尺寸”的步骤}；
放置后可直接拖动文字调整位置。

\subsection{格式}

增减小数位、加 ⌀/□/n× 前缀、去掉前缀、自定义 printf 格式、HoleShaftFit（H7/g6 等配合）、
\ui{修复尺寸}（把尺寸重新指向新的参考）。

\section{装饰}

\begin{itemize}
  \item 中心线（面/两条线/两个点）、修饰顶点/中点/象限点、修饰线/圆/圆弧、
        平行/垂直修饰线、橡皮擦、\ui{显示全部}；
  \item 线属性（线型 1--5、线宽 0.18--0.7）；
  \item 剖面线（Hatch，用 SVG 图案）与几何剖面线（Geometric Hatch，11 种 PAT 图案）。
\end{itemize}

\begin{warn}[装饰几何在页面上的显示有限]
当前版本里，通过“修饰”命令添加的中心线、修饰线等\textbf{不会绘制在页面上}；剖面线统一按斜线绘制，
线型/线宽也不生效。出图时请以视图轮廓 + 尺寸为主，装饰作为补充。
\end{warn}

\section{导出与打印}

\begin{itemize}
  \item \ui{导出为 SVG}：模板 + 轮廓 + 尺寸 + 文字 + 气球（引线、符号、图片、剖面线暂缺）；
  \item \ui{导出为 DXF}：二维 CAD 通用格式；
  \item \ui{导出为 PDF}：按页面 mm 换算成 72 dpi 的矢量绘制；入口在 \ui{更多工具} 里的
        \ui{将页面导出为 PDF}；
  \item \ui{Print All} 与 PDF 导出同一路径：\textbf{只导出当前页}，不经过系统打印对话框。
\end{itemize}

\section{页面交互}

\begin{itemize}
  \item \textbf{平移}：单指在页面空白处拖动；\textbf{缩放}：双指捏合（以两指中点为锚点）；右上角
        \ui{适合} 让整页适配屏幕；切页时自动适合；
  \item \textbf{移动}：拖动已选中的视图/尺寸/气球（视图改 X/Y，尺寸改文字位置，气球改气泡位置），
        一次拖动是一步可撤销操作；
  \item \textbf{选择优先级}：顶点 12 dp → 边 10 dp → 尺寸/气球锚点 18 dp → 面内部（最小面优先）→ 视图边框；
        要点选整个视图，点视图内空白处或在树里点；
  \item \textbf{高亮}：视图边深蓝加粗、顶点蓝点、面半透明蓝、视图虚线框；预选顶点青色；
  \item 提示语：“请先选择视图（或其边/顶点）”“请先在页面上创建页面”等。
\end{itemize}

\begin{practice}[练习：三视图 + 标注 + PDF]
\begin{enumerate}
  \item 打开第 4 章的支架（先保存）；
  \item 切 \ui{工程图} → \ui{新建工程图页}；
  \item \ui{投影组} → 选支架 → 投影选 \ui{前/上/右} → 生成三视图；
  \item 点视图边（变蓝加粗）→ \ui{长度}/\ui{水平}/\ui{垂直} 标外形尺寸；圆边 → \ui{直径}；
        两边 → \ui{角度}；拖尺寸文字调位置；
  \item \ui{更多工具 → 将页面导出为 PDF} → 选保存位置；
  \item 打开 PDF 检查：图纸、视图、尺寸都在。
\end{enumerate}
\end{practice}

\begin{practice}[练习：剖视与局部放大]
\begin{enumerate}
  \item 选前视图 → \ui{剖视图} → 方向“右” → 生成 A-A 剖视；
  \item 选中剖视，点内部面 → \ui{几何剖面线} → 选图案（如 Steel）；
  \item 选基视图 → \ui{局部详图} → 点放大中心 → 半径 10 → 得到 2× 的 Detail A；
  \item 在详图上加直径、中心线。
\end{enumerate}
\end{practice}

\section{源码可见的限制}

\begin{itemize}
  \item 打印只导出当前页，且没有系统打印对话框；
  \item SVG 导出缺引线、符号、图片、剖面线、修饰几何；
  \item 页面不绘制中心线/修饰几何/螺纹/装饰线，线型线宽不生效，剖面线统一斜线；
  \item 断裂视图没有断裂定义；裁剪组不裁剪；图像固定 100×100 mm；
  \item 剖视只能过视图中心四个方向、没有深度参数；局部详图固定 2× 圆形；
  \item 投影组没有第一角/第三角切换；视图没有“当前视图方向”或任意方向；
  \item 尺寸不能点放置；轴向长度退化为长度；弧长以线性尺寸表示；面积/弧长标注画成矩形气球；
  \item 页面比例/幅面只能通过模板或属性修改；“锁定”不阻止拖动。
\end{itemize}
